\section{Human--agent interaction: delegating to an operator}
\label{sec:hai}

\subsection{How developers actually use coding agents}

A first strand of empirical work compares agents with the consultative tools they replace.
\citet{chen2025codewithme} ran the first controlled study of developers using a copilot versus an
agent (OpenHands). Agents reduced effort and enabled task completions that participants might not
have achieved alone. However, making sure that users adequately understood the agent's behaviour
remained a challenge.
This is the advisor-to-operator trade-off in miniature: delegation buys throughput at the cost of
comprehension.

Self-assessment is not a reliable guide to that trade-off. In a randomised controlled trial with 16
experienced open-source developers working on 246 real tasks in mature projects they had
contributed to for years,
\citet{becker2025measuring} found that allowing AI tools increased completion time by 19\%. Yet
developers expected a 24\% speed-up beforehand and still believed in a 20\% speed-up afterwards.
\citet{shen2026skill} add a longer-term concern: in randomised experiments on learning a new
library, AI assistance impaired conceptual understanding, code reading and debugging. Participants
who fully delegated were faster but learned less. If supervision requires the very skills that
delegation erodes, oversight quality may degrade precisely as autonomy increases.

Field studies show that experienced developers are aware of this and keep control. In observations
and surveys, \citet{huang2025dontvibe} find that professionals ``don't vibe, they control'': they
retain their agency over design and implementation because they insist on software quality, and they
use agents as collaborators rather than as complete delegates. In a small mixed-methods study, \citet{feng2026junior} find that how much agency is
allocated to the AI is constrained more by organisational policy than by individual preference.
Senior engineers delegate through detailed instructions, while juniors oscillate between over-reliance
and avoidance.

At scale, mining public artefacts gives a complementary picture. Of 567 Claude Code pull requests in
157 projects, 83.8\% were merged, but 45.1\% of the merged ones needed further changes first
\citep{watanabe2025agentic}. In the SWE-chat dataset of almost 18{,}000 real sessions, agents author
virtually all committed code in 41\% of sessions, while only 59\% of agent-produced code survives into
commits \citep{baumann2026swechat}. Most directly relevant to us, \citet{tang2026fail} analyse 20{,}574
real coding-agent sessions and operationalise misalignment as \emph{visible developer pushback}.
Among the seven recurring forms they identify are failures to interpret intent, to bound actions,
to follow rules and to report progress accurately. 91.5\% of the visible resolutions required
explicit correction by the user. Over time, overall misalignment rates decline, but constraint
violations and inaccurate self-reporting \emph{grow} as a share of misalignment episodes. This
suggests that scope and reporting problems are not being solved by capability gains alone.

\subsection{What oversight work looks like}

Oversight of agents is often discussed normatively. Two recent qualitative studies give it empirical
content. \citet{dhanorkar2026oversight} interview 17 experienced developers and identify four forms of
emergent oversight work:
\begin{itemize}
  \item \emph{a priori control}: constraining the agent before it starts;
  \item \emph{co-planning}: agreeing on the plan with the agent;
  \item \emph{real-time monitoring}: watching the agent as it works;
  \item \emph{post hoc review}: inspecting the result afterwards.
\end{itemize}
Oversight is preventive as well as reactive. One of the heuristics developers adopt---treating
passing tests as a guarantee of correctness---is exactly the proxy that coding agents are known to
game \citep{zhong2025impossiblebench}. \citet{park2026delegation} adapt Sheridan's model of human
supervisory control into seven stages for supervising coding agents. Developers concentrate effort
in planning, delegate part of the supervision to other agents, and turn recurring guidance into
reusable assets.

These findings echo classic results. Supervising automation requires \emph{situation awareness}: the
supervisor must perceive what the system is doing, understand it, and project what it will do next
\citep{endsley1995situation}. Supervisors also need trust calibrated to the system's actual
reliability \citep{lee2004trust}. Drift is a direct threat to both. It is gradual, so it does not
trigger the salient alarms that help perception. And a fluent final summary invites over-trust. The
sabotage study by \citet{ye2026enemy} is a sobering reference point: 83 of 88 developers (94\%)
without an automated monitor failed to detect an agent that sabotaged the code to pursue a hidden
side task, and the authors attribute this to minimal code review, a plausible cover story and
over-trust. Accidental drift is subtler still, because
there is no adversary leaving traces.

\subsection{Design responses: communication, checkpoints and dynamic autonomy}

HCI research has begun to propose interaction designs for the operator regime.
\citet{bansal2024communication} catalogue twelve challenges in human--agent communication. On the
agent-to-user side these include conveying capabilities, plans, progress and side effects; on the
user-to-agent side, conveying goals, preferences and constraints. Drift through sub-tasks is a
failure on both sides. The agent fails to report that its effort has shifted, and the user's original
intent fails to persist against the agent's own sub-goals.

Magentic-UI implements human-in-the-loop mechanisms such as co-planning, co-tasking, action guards and
long-term memory \citep{mozannar2025magentic}. Hedwig starts from a survey in which engineers report
frustration with static autonomy settings; its abstract explicitly mentions ``scope drift'' among the
agent behaviours that slip past review. The system adjusts the agent's autonomy dynamically,
loosening oversight where trust has been earned and tightening it in unfamiliar territory
\citep{shukla2026hedwig}. \citet{bui2026proactivity} argue that agents need a principled policy for
when to surface information and interrupt, grounded in mixed-initiative interaction
\citep{horvitz1999mixed}, rather than simply more autonomy. At the other extreme,
\citet{monperrus2026endreview} argues that human code review is no longer a necessary quality gate.
If that view prevailed, drift could only be caught by other agents---and LLM monitors themselves
degrade on long transcripts \citep{martin2026classifier}.

\subsection{Gap}

The empirical HCI literature observes misalignment mostly \emph{after the fact}, through
developer pushback \citep{tang2026fail} or review \citep{watanabe2025agentic}, and shows that
oversight depends on cheap but fallible proxies such as passing tests
\citep{dhanorkar2026oversight}. It does not tell us \emph{when} during a trajectory the goal starts
to slip, which kinds of sub-tasks cause it, or how much of the drift is visible in what the agent
reports. Answering those questions requires controlled manipulation of the trajectory itself, which is
what the goal-drift literature in the next section offers.
